\chapter{Podsumowanie}\label{chapter:conclusions}

W trakcie realizacji pracy inżynierskiej został zaprojektowany i zaimplementowany system zarządzania projektami ITForge, który stanowi kompleksowe narzędzie wspierające proces zarządzania projektami informatycznymi. Niniejszy rozdział zawiera podsumowanie wykonanych prac, ocenę realizacji celów projektu, wnioski wynikające z przeprowadzonych prac oraz możliwości dalszego rozwoju systemu.

\section{Wykonane prace}

W ramach realizacji pracy inżynierskiej wykonano szereg prac obejmujących wszystkie etapy cyklu życia oprogramowania, od analizy wymagań po implementację i testowanie systemu.

W fazie analizy i projektowania przeprowadzono szczegółową analizę wymagań funkcjonalnych i niefunkcjonalnych systemu, zdefiniowano przypadki użycia dla wszystkich głównych modułów, zaprojektowano architekturę systemu opartą na wzorcu klient-serwer oraz opracowano model bazy danych obejmujący tabele dla użytkowników, projektów, zadań, wymagań, ryzyk, backlogów i sprintów. Przygotowano diagramy przypadków użycia oraz diagramy czynności dla kluczowych procesów biznesowych, takich jak dodawanie zadań w tablicy Kanban i wykresie Gantta.

W fazie implementacji stworzono pełnofunkcjonalną aplikację webową opartą na architekturze klient-serwer. Warstwa frontendowa została zrealizowana przy użyciu biblioteki React wraz z React Router do zarządzania routingiem oraz Context API do zarządzania stanem globalnym aplikacji, tworząc responsywny i interaktywny interfejs użytkownika. Zaimplementowano komponenty dla wszystkich głównych modułów systemu, w tym panelu projektów, tablicy Kanban, wykresu Gantta, zarządzania wymaganiami, zarządzania ryzykiem, backlogów i sprintów oraz zarządzania zespołem. Warstwa backendowa została zrealizowana przy użyciu Node.js i frameworka Express.js, tworząc RESTful API obsługujące wszystkie funkcjonalności systemu. Zaimplementowano mechanizmy autentykacji i autoryzacji oparte na tokenach JWT oraz hashowaniu haseł przy użyciu biblioteki bcrypt, zapewniając bezpieczeństwo danych użytkowników. Baza danych została zrealizowana w systemie PostgreSQL, z zastosowaniem migracji zapewniających kontrolę wersji schematu bazy danych.

Zrealizowano następujące moduły funkcjonalne: system autentykacji i autoryzacji użytkowników z możliwością rejestracji i logowania, moduł zarządzania projektami umożliwiający tworzenie, edycję, usuwanie oraz przeglądanie projektów, system zarządzania zespołem z mechanizmem zaproszeń i powiadomień, tablicę Kanban z możliwością tworzenia, edycji i usuwania kolumn oraz zadań, z funkcją przeciągania i upuszczania zadań między kolumnami, wykres Gantta umożliwiający wizualizację zadań w czasie z możliwością definiowania zależności między zadaniami, moduł zarządzania wymaganiami projektowymi z możliwością definiowania, edycji i usuwania wymagań, moduł zarządzania ryzykiem z możliwością identyfikacji, oceny i monitorowania ryzyk oraz edycji macierzy ryzyka, oraz moduł backlogów i sprintów umożliwiający pracę w metodyce Scrum z możliwością tworzenia sprintów i zarządzania elementami backlogu.

W fazie testowania przeprowadzono kompleksowe testy funkcjonalne obejmujące 75 przypadków testowych, które pokryły wszystkie zdefiniowane w założeniach projektowych funkcjonalności. Testy obejmowały autentykację i autoryzację użytkowników, zarządzanie projektami i członkami zespołu, funkcjonalności tablicy Kanban, funkcjonalności wykresu Gantta, zarządzanie wymaganiami projektowymi, zarządzanie backlogami i sprintami oraz zarządzanie ryzykiem projektowym. Wszystkie testy zakończyły się wynikiem pozytywnym, co potwierdza poprawność implementacji systemu i spełnienie założonych wymagań funkcjonalnych.

\section{Ocena realizacji celów projektu}

Głównym celem pracy było stworzenie funkcjonalnego systemu zarządzania projektami, który integruje różnorodne metodyki i narzędzia wspierające proces zarządzania projektami. W wyniku przeprowadzonych prac udało się zrealizować wszystkie założone cele projektu.

Pierwszym zrealizowanym celem było stworzenie systemu umożliwiającego kompleksowe zarządzanie projektami, w tym tworzenie, edycję i usuwanie projektów, zarządzanie zespołem oraz kontrolę dostępu do zasobów projektowych. System w pełni spełnia te wymagania, oferując intuicyjny interfejs do zarządzania projektami oraz mechanizmy kontroli dostępu oparte na rolach użytkowników.

Drugim zrealizowanym celem była integracja metod zarządzania zadaniami poprzez zaimplementowanie dwóch komplementarnych metod wizualizacji: tablicy Kanban oraz wykresu Gantta. Oba moduły działają poprawnie i pozwalają na elastyczne dostosowanie narzędzi do preferencji użytkowników i charakteru projektu. Tablica Kanban umożliwia wizualne zarządzanie przepływem pracy, podczas gdy wykres Gantta wspiera planowanie czasowe i zarządzanie zależnościami między zadaniami.

Trzecim zrealizowanym celem było utworzenie systemu umożliwiającego definiowanie, śledzenie i zarządzanie wymaganiami projektowymi. Moduł wymagań został w pełni zaimplementowany i pozwala na kompleksowe zarządzanie wymaganiami, co jest kluczowe dla zapewnienia jakości realizowanych projektów.

Czwartym zrealizowanym celem było zaimplementowanie modułu pozwalającego na identyfikację, ocenę i monitorowanie ryzyk projektowych wraz z macierzą ryzyka. Moduł zarządzania ryzykiem został w pełni zrealizowany i wspiera proces podejmowania decyzji zarządczych poprzez wizualizację ryzyk w macierzy ryzyka.

Piątym zrealizowanym celem było stworzenie modułu backlogów i sprintów, umożliwiającego pracę w metodyce Scrum. System zawiera pełnofunkcjonalny moduł wspierający metodykę Scrum, co czyni go przydatnym dla zespołów zwinnych.

Szóstym zrealizowanym celem było zaimplementowanie mechanizmu powiadomień oraz zaproszeń do projektów. System zawiera funkcjonalny moduł powiadomień i zaproszeń, co ułatwia komunikację i współpracę w zespole.

Realizacja projektu przebiegła zgodnie z założonym harmonogramem i pozwoliła na osiągnięcie wszystkich głównych celów. System ITForge stanowi funkcjonalne narzędzie wspierające proces zarządzania projektami informatycznymi, integrujące różnorodne metodyki i narzędzia w jednej platformie. Projekt pozwolił na praktyczne zastosowanie wiedzy z zakresu inżynierii oprogramowania, baz danych, technologii webowych oraz zarządzania projektami. Zrealizowany system może stanowić podstawę do dalszego rozwoju i wdrożenia w rzeczywistych warunkach.

\section{Wnioski}

W zakresie aspektów technicznych zastosowanie architektury klient-serwer z podziałem na warstwę prezentacji (React), warstwę logiki biznesowej (Node.js/Express) oraz warstwę danych (PostgreSQL) okazało się właściwym wyborem, zapewniając dobrą separację odpowiedzialności i możliwość niezależnego rozwoju poszczególnych komponentów. Wykorzystanie biblioteki React wraz z React Router do zarządzania routingiem oraz Context API do zarządzania stanem globalnym aplikacji pozwoliło na stworzenie responsywnego i interaktywnego interfejsu użytkownika. Zastosowanie komponentów funkcyjnych z hookami React ułatwiło zarządzanie stanem komponentów i logiką biznesową. Framework Express.js okazał się wystarczający do zrealizowania wszystkich wymaganych funkcjonalności. Zastosowanie middleware do autentykacji i autoryzacji zapewniło bezpieczeństwo systemu. Struktura oparta na modelach i kontrolerach ułatwiła utrzymanie czytelności kodu. Wykorzystanie systemu zarządzania bazą danych PostgreSQL pozwoliło na efektywne przechowywanie i zarządzanie danymi. Zastosowanie migracji bazy danych ułatwiło zarządzanie schematem bazy danych i zapewniło kontrolę wersji struktury danych. Zaimplementowano mechanizmy autentykacji oparte na tokenach JWT oraz hashowanie haseł użytkowników przy użyciu biblioteki bcrypt, co zapewnia odpowiedni poziom bezpieczeństwa danych użytkowników.

W trakcie realizacji projektu napotkano na szereg wyzwań technicznych i projektowych. Jednym z głównych wyzwań było zapewnienie spójności danych między różnymi modułami systemu (Kanban, Gantt, Backlog), szczególnie w kontekście zadań, które mogą być wyświetlane w różnych widokach. W początkowych etapach rozwoju systemu niektóre zapytania do bazy danych wymagały optymalizacji, szczególnie w kontekście pobierania danych powiązanych (np. zadania wraz z przypisanymi użytkownikami). Wraz z rozwojem aplikacji zarządzanie stanem globalnym stawało się coraz bardziej złożone. Wprowadzenie Context API pomogło w rozwiązaniu tego problemu, choć w większych aplikacjach warto rozważyć wykorzystanie bardziej zaawansowanych rozwiązań, takich jak Redux. Integracja biblioteki do wizualizacji wykresu Gantta wymagała dostosowania do specyfiki projektu oraz zapewnienia odpowiedniej synchronizacji danych z backendem.

W trakcie realizacji projektu zidentyfikowano następujące ograniczenia obecnej wersji systemu. System nie posiada mechanizmu aktualizacji danych w czasie rzeczywistym (np. WebSockets), co oznacza, że zmiany wprowadzone przez jednego użytkownika nie są natychmiast widoczne dla innych użytkowników bez odświeżenia strony. System nie zawiera zaawansowanych funkcji raportowania i analityki, które mogłyby wspierać proces podejmowania decyzji zarządczych. System nie posiada możliwości integracji z popularnymi narzędziami, takimi jak systemy kontroli wersji (Git), systemy komunikacyjne (Slack, Teams) czy narzędzia do zarządzania dokumentacją. Obecna architektura systemu może wymagać optymalizacji w przypadku obsługi bardzo dużej liczby użytkowników lub projektów jednocześnie.

Warto podkreślić, że system został w pełni przetestowany i wszystkie zdefiniowane funkcjonalności działają poprawnie, co potwierdza wysoką jakość implementacji. Pozytywne wyniki testów funkcjonalnych potwierdzają, że system działa zgodnie z założeniami projektowymi i jest gotowy do użycia. Jednocześnie zidentyfikowane ograniczenia i kierunki rozwoju wskazują na potencjał dalszego doskonalenia systemu.

\section{Możliwości dalszego rozwoju}

Na podstawie przeprowadzonej analizy oraz zidentyfikowanych ograniczeń, można wskazać następujące kierunki dalszego rozwoju systemu.

Wprowadzenie mechanizmu WebSockets lub Server-Sent Events (SSE) umożliwiłoby aktualizację danych w czasie rzeczywistym, co znacznie poprawiłoby doświadczenie użytkownika, szczególnie w kontekście współpracy zespołowej. Umożliwiłoby to natychmiastowe wyświetlanie zmian wprowadzonych przez innych użytkowników bez konieczności odświeżania strony.

Dodanie zaawansowanych funkcji raportowania, w tym wykresów, statystyk oraz eksportu danych do różnych formatów (PDF, Excel), wsparłoby proces analizy i monitorowania postępów projektów. Raporty mogłyby obejmować statystyki dotyczące postępów w realizacji zadań, analizę wykorzystania czasu, raporty z sprintów oraz analizę ryzyk projektowych.

Wprowadzenie API umożliwiającego integrację z popularnymi narzędziami deweloperskimi i komunikacyjnymi zwiększyłoby użyteczność systemu i ułatwiło jego wdrożenie w istniejących środowiskach pracy. Integracje mogłyby obejmować systemy kontroli wersji (Git, GitHub, GitLab), systemy komunikacyjne (Slack, Microsoft Teams), narzędzia do zarządzania dokumentacją (Confluence, Notion) oraz narzędzia do ciągłej integracji i wdrażania (CI/CD).

Dodanie możliwości przypisywania etykiet, priorytetów, szacowania czasu oraz śledzenia czasu pracy nad zadaniami wzbogaciłoby funkcjonalność systemu. Umożliwiłoby to bardziej szczegółowe zarządzanie zadaniami oraz lepsze planowanie i śledzenie postępów w realizacji projektów.

Rozszerzenie systemu powiadomień o powiadomienia e-mail oraz możliwość konfiguracji preferencji powiadomień zwiększyłoby wygodę użytkowania. Użytkownicy mogliby otrzymywać powiadomienia o ważnych wydarzeniach w projektach, takich jak nowe zaproszenia, zmiany w zadaniach czy zbliżające się terminy.

Rozwój aplikacji mobilnej (iOS/Android) lub responsywnej wersji webowej zoptymalizowanej pod urządzenia mobilne umożliwiłby dostęp do systemu z dowolnego miejsca. Aplikacja mobilna pozwoliłaby użytkownikom na zarządzanie projektami i zadaniami w trakcie przemieszczania się, co jest szczególnie istotne w kontekście współczesnych metod pracy zdalnej.

Wprowadzenie dodatkowych mechanizmów bezpieczeństwa, takich jak uwierzytelnianie dwuskładnikowe (2FA), logowanie aktywności użytkowników oraz bardziej zaawansowane mechanizmy kontroli dostępu, zwiększyłoby bezpieczeństwo systemu. Uwierzytelnianie dwuskładnikowe zapewniłoby dodatkową warstwę ochrony kont użytkowników, podczas gdy logowanie aktywności umożliwiłoby śledzenie działań użytkowników w systemie.

Wprowadzenie mechanizmów cache'owania, optymalizacji zapytań do bazy danych oraz implementacja load balancing w przypadku większego obciążenia systemu poprawiłoby wydajność i skalowalność systemu. Mechanizmy te byłyby szczególnie istotne w przypadku wdrożenia systemu w środowisku produkcyjnym z dużą liczbą użytkowników.

Podsumowując, realizacja pracy inżynierskiej pozwoliła na stworzenie kompleksowego systemu zarządzania projektami ITForge, który spełnia założone wymagania funkcjonalne i techniczne. System integruje różnorodne narzędzia i metodyki zarządzania projektami, co czyni go uniwersalnym rozwiązaniem dla zespołów projektowych. Pozytywne wyniki testów funkcjonalnych potwierdzają poprawność implementacji, a zidentyfikowane kierunki rozwoju wskazują na możliwości dalszego doskonalenia systemu w przyszłości.
